\section{Scalar exponents on a compact threefold}\label{sec:compact}

\begin{corollary}\label{cor:arith-roots}
Let the image of $\alpha$ in $H_2(X;\ZZ)/\mathrm{tors}$ be primitive,
where $X$ is a smooth projective Calabi--Yau threefold. For each $g$
let $C_g=\sum_{3\gamma=3\alpha}c_{\gamma,g}$, summing over the
effective integral roots. The common denominator of the scalar product
exponents $e_{3\alpha,h}$ is $3^M$, where
\[
 M=\max\bigl(\{0\}\cup
       \{s_n-v_3(C_{n+1}): n\ge0,\ C_{n+1}\ne0\}\bigr).
\]
In particular all these exponents are integral if and only if
$3^{s_n}\mid C_{n+1}$ for every $n\ge0$.
\end{corollary}
\begin{proof}
By \eqref{eq:compact-prime}, the polynomial of differences at this
class is $\sum_{n\ge0}C_{n+1}D_n$. There are finitely many occupied
genera, and $c_{3\alpha,h}$ is an integer for each $h$ by
\cite[Theorem~1.1]{WueProduct}. For $M>0$ scale by $3^M$ and reduce
modulo three. Each maximizing index contributes a nonzero unit times
the corresponding reduced polynomial of Theorem~\ref{thm:arith-rows}.
Their monomial supports, $\{9k\}$ or $\{9k+2,9k+3\}$, are disjoint
for distinct indices. This proves equality of the denominator; the
sum over torsion-distinct roots of a single genus has already been
taken in $C_g$. For $M=0$ all rows in the sum are integral.
\end{proof}

In particular \eqref{eq:arith-row} gives
\begin{equation}\label{eq:arith-ten}
 e_{3\alpha,10}=c_{3\alpha,10}-\tfrac{4720}{3}C_3.
\end{equation}
Indeed $D_1$ has degree six, and for $n\ge3$ the smallest possible
power of $w$ is $3n+1$ by Theorem~\ref{thm:main}.

Let $Y$ be the smooth anticanonical threefold in the toric resolution
of $X_9$ whose fan is specified in Appendix~\ref{sec:geometry}. Its second homology is
torsion-free and it contains a divisor $F\simeq\mathbb F_1$ with
normal line $K_F$. Let $s,f$ be the negative section and a ruling fiber.
For $a=12j+3$, $j\ge0$, set
\[
 \alpha_a=as+(a+1)f,\quad G=a(a-1)/2,\quad k=(G-1)/2.
\]
We use the same notation for curve classes on $F$ and their images in
$Y$. Two integral nef divisors on $Y$ recover the coefficients of
$s,f$ by intersection. Their difference takes value one on $\alpha_a$.
The four other nef degrees of \eqref{eq:degrees} vanish on both
$s$ and $f$. Their common zero face confines every effective curve of
class $\alpha_a$, and every effective summand, to $F$. Thus $\alpha_a$
has primitive free image and it is the sole integral root at
$3\alpha_a$.

These are the classes in Theorem~\ref{thm:arith-compact}.

\subsection{The highest primitive genus}

Put $D=as+(a+1)f$. In the presentation $F=\operatorname{Bl}_p\PP^2$
one has $D=(a+1)H-E$. Its arithmetic genus is $G$, and
\[
 N=h^0(F,\mathcal O_F(D))=\binom{a+3}{2}-1
                          =\tfrac12(a+1)(a+4).
\]
We show that the compact elementary exponents satisfy
\begin{equation}\label{eq:arith-minimum}
 c_{\alpha_a,h}=0\quad(h>G),\qquad
 c_{\alpha_a,G}=(-1)^{N-1}N.
\end{equation}

First consider stable pairs on $Z=\operatorname{Tot}(K_F)$. All
positive compact curves are supported on $F$: affine fibers contain
no complete curves, and on a normalization projecting nontrivially
to $F$ the tautological section of the negative canonical line
vanishes. For an effective decomposition $D=\sum D_i$, with at
least two nonzero classes $D_i=a_i s+b_i f$, $a_i,b_i\ge0$, we claim
\begin{equation}\label{eq:arith-split}
 \sum_i\chi(\mathcal O_{D_i})-\chi(\mathcal O_D)
 =\sum_{i<j}D_iD_j
 =\sum_i b_i(a-a_i)-\sum_{i<j}a_i a_j>0.
\end{equation}
These numbers depend only on the divisor classes by Riemann--Roch.
Choose $A=\max a_i$ and $t=a-A$. For $t>0$ the first sum in the
last expression is at least $(a+1)t$, and the second is at most
$At+t(t-1)/2$. Their difference is at least $t(t+3)/2$. For $t=0$
some positive multiple of a fiber lies in a different summand, giving
an intersection of at least $a$.

For a pure curve $W\subset Z$ filter $\mathcal O_W$ by powers of
the ideal of $F$. The graded quotients are of the form
$K_F^{-i}\otimes\mathcal O_F/J_i$. The divisorial hull of $J_i$ is
$\mathcal O_F(-D_i)$, and its quotient by $J_i$ has finite length
$l_i$. The sum of the divisor classes $[D_i]$ is $[W]$. The Euler
characteristic of the $i$th layer is
$\chi(\mathcal O_{D_i})-iK_FD_i+l_i$, with both corrections
nonnegative. By \eqref{eq:arith-split}, two nonzero divisor layers
in class $D$ give $\chi(\mathcal O_W)>1-G$. If there is only one,
it is the zeroth layer, since a nonzero higher layer along a curve
requires a nonzero zeroth layer along that curve. Then
$I_F\mathcal O_W$ has finite support and vanishes by purity. Such a
curve is an effective Cartier divisor in $|D|$. The section of a
stable pair has image $\mathcal O_W$ and finite-length cokernel.
It follows that the minimum pair Euler characteristic is $1-G$ and
all pairs at the minimum are quotient pairs in $|D|$. For a product
of pairs in a nontrivial additive splitting of $D$, apply
\eqref{eq:arith-split} to their combined divisor layers. Its total
Euler characteristic is strictly greater than $1-G$.

The minimum pair space is $\PP^{N-1}$ as a scheme. Indeed
surjectivity of the section in a family follows from Nakayama's
lemma, identifying it with the corresponding Hilbert functor.
For every $W\in|D|$, the normal sequence for $W\subset F\subset Z$
has terms $\mathcal O_W(D)$, $N_{W/Z}$ and $K_F|_W$. The last has
no sections since $H^0(F,K_F)=0$ and $H^1(F,K_F-D)=0$. The latter
follows from duality and Kodaira vanishing, as $D-K_F$ is ample.
Thus $h^0(N_{W/Z})=h^0(\mathcal O_W(D))=N-1$ by rationality of
$F$. The closed immersion of $|D|$ exhausts the closed points, so the
local dimension and the embedding dimension of the Hilbert space
both equal $N-1$ there. Its local rings are regular, proving the
scheme assertion. The symmetric pair obstruction theory on this smooth
proper scheme has degree $(-1)^{N-1}N$ \cite{PT}. The strict product
inequality implies that the coefficient of the formal logarithm of
the stable-pair series at $Q^D$ has the same lowest term,
$(-1)^{N-1}N z^{1-G}$ in the pair variable $z$.

To identify this term with a compact elementary exponent, deform $Y$
to the normal line of $F$. Blow up $Y\times\mathbb A^1$ at
$F\times0$ and remove the strict transform of $Y\times0$. The resulting
smooth family has fibers $Y$ off zero and $Z$ over zero and contains
a Cartier divisor $F\times\mathbb A^1$. In each positive face class
all connected maps have image in that divisor by the nef degrees on
$Y$ and negativity on $Z$. This is true as a statement about families
of maps: on every geometric source the Cartier line restricts to the
pullback of $K_F$. Sections vanish on every positive component by
negativity, then on all contracted components by connectedness. In a
two-term presentation of the derived direct image the first differential
is therefore fiberwise injective and locally split injective. The
canonical section vanishes on the universal domain, including over
nonreduced bases. Hence the relative moduli is the product of the
proper map stack of $F$ and the parameter line. Deformation invariance
of its ordinary relative stable-map class gives equality of the
connected GW degrees of $Y$ and $Z$ in each face class.

The fixed-target toric correspondence \cite[Theorem~1]{PP} applies
to $Z$. The pair spaces in these classes are proper: in the projective
line completion of $K_F$ their geometric supports are on the zero
section, and a flat specialization of a pair supported there remains
set-theoretically on that section. For instance a power of its ideal
annihilating the generic pair sheaf annihilates the flat family over
a discrete valuation ring. Thus the equivariant degrees have their
ordinary specialization. Since $K_Z$ is trivial there is no prefactor
in the empty-insertion correspondence under $z=-e^{iu}$. Taking the
formal logarithm through class $D$ uses only the finitely many
positive summands on this face, with disconnected maps nonconstant on
each connected component. The result identifies the connected fixed-target
GW series for $Y$ in $\alpha_a$ with the rational function of $z$
whose lowest term was computed above.

On the compact projective $Y$, the algebraic and symplectic GW degrees
agree \cite{Siebert}. At a class with primitive free image the finite
elementary-cluster expansion of \cite{IP,DIW} is
$\sum_h c_{\alpha_a,h}w^{h-1}$, with these same compact exponents by
\cite[Theorem~1.1]{WueProduct}. Since $w=(1+z)^2/z$, this is a rational
function of $z$. Equality of the expansions at $z=-1$ gives equality
of rational functions. At $z=0$ the summand of largest genus $h$ has
lowest Laurent power $1-h$ with coefficient $c_{\alpha_a,h}$.
Comparison with the lowest pair term proves~\eqref{eq:arith-minimum}.

\begin{proof}[Proof of Theorem~\ref{thm:arith-compact}]
For $a=12j+3$ the integer $G=2k+1$ is odd and $N\equiv2\pmod3$.
In \eqref{eq:compact-prime} the top occupied row is $D_{2k}$.
Theorem~\ref{thm:arith-rows} shows that its denominator is $3^k$ and
that $3^kD_{2k}$ reduces to $2(-1)^kw^{9k}$. Every lower row has
denominator dividing $3^{k-1}$. The coefficient multiplying the top
row is a unit at three by~\eqref{eq:arith-minimum}, while the
compact elementary exponents are integral. This proves the exact
denominator at output genus $9k+1$. For $a=3$ one has $G=3$, $N=14$
and $c_{\alpha_3,3}=-14$. Equation~\eqref{eq:arith-ten} gives
$e_{3\alpha_3,10}=c_{3\alpha_3,10}+66080/3$. For $a=15$ the
corresponding genus is $469$ and the denominator is $3^{52}$.
\end{proof}

\subsection{The first possible fractional exponents}
If the free image of a class is primitive its scalar and opposite-weight
exponents agree. At free divisibility two, \eqref{eq:compact-prime}
and the integral coefficients of \eqref{eq:arith-two} imply
integrality in all genera, summing over every integral double root.
There is also a uniform statement in small genus. For all effective
$\beta$ on a smooth projective Calabi--Yau threefold as above,
\[
 e_{\beta,h}=c_{\beta,h}\quad(0\le h\le3),\qquad
 e_{\beta,4}=c_{\beta,4}-2\sum_{2\gamma=\beta}c_{\gamma,2}.
\]
Indeed pull the logarithms of the two compact products back to
connected maps in source genus at most four. A positive-degree
map to a core of genus $g\ge2$ in projection degree $d$ has source
genus at least $1+d(g-1)$ by normalization and Riemann--Hurwitz.
The zero and one core genera agree in every degree. At the remaining
small degrees Theorem~\ref{thm:main} and \eqref{eq:arith-two} leave
precisely the term $2x_{4,1}$ from $(g,d)=(2,2)$ through source genus
four. This assertion also holds for the equivariant cycle insertions:
their supported pullbacks are over the zero-section multiple, hence
empty below that same genus. Induct on ample degree, reading source
genera in increasing order by the proper detection of
\cite[Sections~5--6]{WueProduct}. Proper multiples of a genus-four change
can first be detected in genus seven. This gives the displayed identities.

One may also try to deduce integrality from the scalar stable-pair
series. If all scalar local pair coefficients were integral Laurent
series in the pair variable $z$, division of a compact integral pair
series by already determined lower factors would leave an integral
coefficient $\sum_g e_{\beta,g}z^{1-g}(1+z)^{2g-2}$ in the next
class. The largest occupied genus has a unique lowest Laurent term
with unit coefficient; descending induction would give integral
exponents. The local hypothesis fails. In the genus-three scalar
cubic the disconnected trace in \eqref{eq:scalar-traces}, under
$w=(1+z)^2/z$, has coefficient $20+292/3=352/3$ at $z^{-9}$.
Here one uses the fixed-target equivariant local GW/pairs correspondence
\cite[Section~4.2]{Pardon}; the lowest compact proper pair count used
in Theorem~\ref{thm:arith-compact} is integral.
